\noindent
\ProjectTitle{KōreroCare: Secure Mobile–Wearable Digital Phenotyping Infrastructure (2026 | Research)}

\noindent
\textit{Associated with Empathic Computing Laboratory}

\begin{itemize}
    \item Designed a research-grade digital phenotyping infrastructure for longitudinal behavioral and physiological data collection across Android, WearOS, and backend services.

    \item Engineered reliable passive sensing workflows using Foreground Services and WorkManager to sustain high-uptime data capture under Android OS constraints.

    \item Built resilient synchronization, local buffering, and retry mechanisms to reduce data loss and improve consistency across mobile and wearable data streams.

    \item Contributed to backend architecture using Go, with service-oriented design for scalable ingestion and secure processing of high-volume time-series data.

    \item Integrated Keycloak for authentication and identity management, supporting token-based access control and extensible researcher/user authentication flows.

    \item Introduced Traefik as an API gateway and load-balancing layer for centralized routing and secure service exposure in a containerized environment.

    \item Explored Open Wearables integration to support interoperable ingestion and normalization of wearable data across heterogeneous device ecosystems.

    \item Defined structured schemas for raw sensor streams and derived features, supporting downstream analytics, machine learning, and longitudinal digital phenotype analysis.

    \item Produced architecture documentation and monitoring-focused designs to improve reproducibility, maintainability, and real-world deployment readiness.
\end{itemize}

\noindent
\textit{Technologies:} Kotlin, Android, WearOS, Go, Keycloak, Traefik, Open Wearables, FastAPI, Docker, WorkManager, Foreground Services
